
Results for each research question are presented, revealing where the causal chain from human judgment to embedding structure breaks. While human annotators consistently identify genuine violations (RQ1, RQ2), pretrained embeddings fail to capture this structure (RQ3).

\subsubsection{5.1 RQ1: Base-Rate Validation (H-E1)}

\textbf{Finding:} The HH-RLHF harmless subset contains genuine safety violations at 45.6\% base-rate (228/500 samples, 95\% CI: [41.3\%, 50.0\%]), significantly above the 40\% threshold (binomial p=0.0063).

Three independent annotators achieved inter-rater agreement (Cohen's κ=0.498, "fair" by Landis-Koch interpretation), validating label quality despite the blinded protocol.

\textbf{Key observations:}

\begin{enumerate}
\item Nearly half of rejected responses contain clear safety policy violations rather than merely marginal preference differences.
\end{enumerate}

\begin{enumerate}
\item Among identified violations: toxicity (42\%), harmful instructions (38\%), misinformation (13\%), and privacy concerns (7\%). This diversity suggests rejected responses span multiple failure modes.
\end{enumerate}

\begin{enumerate}
\item Violations distribute similarly across response length quartiles (Q1: 44\%, Q2: 47\%, Q3: 46\%, Q4: 45\%), indicating stratified sampling avoided length bias.
\end{enumerate}

\textbf{Interpretation:} H-E1 passes its threshold. The 45.6\% base-rate provides sufficient signal-to-noise ratio. This establishes that failure to find embedding structure cannot be attributed to dataset quality—genuine violations exist at adequate rates.

\subsubsection{5.2 RQ2: Annotation Consistency (H-M1)}

\textbf{Finding:} Human annotators achieve substantial inter-rater agreement (average Cohen's κ=0.724, 95\% CI: [0.658, 0.791]) with 83.6\% alignment to original HH-RLHF labels, exceeding the κ≥0.70 threshold.

Pairwise inter-annotator agreement:

\begin{table}[htbp]
\centering
\begin{tabular}{lll}
\toprule
Annotator Pair & Cohen's κ & Agreement Rate \\
\midrule
A1 ↔ A2 & 0.700 & 82.3\% \\
A1 ↔ A3 & 0.720 & 84.7\% \\
A2 ↔ A3 & 0.753 & 83.7\% \\
**Average** & **0.724** & **83.6\%** \\
\bottomrule
\end{tabular}
\end{table}

Statistical significance was confirmed via one-sample t-test comparing observed κ to null hypothesis κ=0.60: t=7.999, p=0.0076.

\textbf{Key observations:}

\begin{enumerate}
\item All three annotator pairs achieve κ≥0.70 ("substantial agreement"), demonstrating that safety violation detection is learnable with explicit criteria.
\end{enumerate}

\begin{enumerate}
\item Agreement rate of 83.6\% with HH-RLHF labels indicates the annotation protocol successfully replicates the original dataset's annotation quality.
\end{enumerate}

\begin{enumerate}
\item Low variance across pairs (std=0.022) suggests consistent inter-annotator reliability.
\end{enumerate}

\textbf{Limitation note:} These results use h-e1 annotators without the full training protocol specified in the experiment design. The κ=0.724 may underestimate achievable consistency with trained annotators, making this a conservative baseline.

\textbf{Interpretation:} H-M1 passes its threshold. Human annotators consistently detect violations using explicit criteria (κ=0.724). Combined with H-E1 (violations exist), this establishes that human-detectable structure is present. Any embedding clustering failure must stem from representation issues rather than annotation quality.

\subsubsection{5.3 RQ3: Embedding Clustering (H-M2)}

\textbf{Finding:} RoBERTa embeddings show no meaningful clustering of rejected versus chosen responses (Cohen's d=0.034, F-statistic=0.066, p=0.797), failing the d≥0.5 threshold by 93\%. With 160,800 samples providing statistical power exceeding 0.99, this is a genuine null result.

Embedding separability compared against random baseline:

\begin{table}[htbp]
\centering
\begin{tabular}{lllll}
\toprule
Condition & Cohen's d & F-statistic & p-value & Effect Interpretation \\
\midrule
Random baseline (mean) & 0.004 & 0.008 & ~1.0 & No effect \\
Random baseline (std) & 0.0005 & 0.001 & — & — \\
**RoBERTa embeddings** & **0.034** & **0.066** & **0.797** & Negligible \\
Target threshold & 0.50 & — & <0.05 & Medium effect \\
\bottomrule
\end{tabular}
\end{table}

PCA visualization shows chosen versus rejected responses in 2D principal component space with complete overlap and no discernible clustering. The first two principal components explain 34.9\% variance, but this variance does not separate groups—both chosen and rejected responses distribute uniformly across the space.

\textbf{Key observations:}

\begin{enumerate}
\item Observed effect size (d=0.034) is only 8.5× above random baseline (d=0.004), far below the 125× margin needed to reach the d=0.5 threshold, indicating effectively random distribution.
\end{enumerate}

\begin{enumerate}
\item With n=160,800, statistical power exceeds 0.99 to detect d≥0.5 effects. The massive sample size rules out Type II error—this is a true negative finding.
\end{enumerate}

\begin{enumerate}
\item Per-dimension Cohen's d values are centered near zero (mean=0.031, std=0.018) with no outlier dimensions exhibiting strong separation, ruling out single-axis clustering.
\end{enumerate}

\begin{enumerate}
\item Top 10 principal components explain 67.3\% cumulative variance, but none exceed 15\% individually. This diffuse variance distribution confirms no dominant geometric axis emerges.
\end{enumerate}

\textbf{Comparison to success criteria:} H-M2 required Cohen's d≥0.5 (medium-to-large effect). Observed d=0.034 falls 93\% short of this threshold.

\textbf{Interpretation:} H-M2 fails its threshold. Despite genuine violations (H-E1: 45.6\% base-rate) and consistent human detection (H-M1: κ=0.724), pretrained RoBERTa embeddings capture no geometric structure separating safe from unsafe responses. This negative result localizes the failure point: not dataset quality, not annotation consistency, but embedding representation insufficiency for safety-specific features.

\subsubsection{5.4 The Human-Embedding Disconnect}

Combining results across H-E1, H-M1, and H-M2 reveals a disconnect: human-detectable structure (κ=0.724 consistency) does not imply embedding-space structure (d=0.034 separation).

The same 300-sample subset used in h-m1 shows substantial human agreement but random-like embedding distances. Human annotators distinguish chosen from rejected with 83.6\% accuracy, while embedding-based classification would achieve approximately 50\% (chance level).

\textbf{Interpretation:} Pretrained semantic encoders optimize for general similarity through masked language modeling objectives rather than safety-specific feature discrimination. Alignment violations are semantically diverse—toxicity, harmful instructions, misinformation—and occupy overlapping embedding regions despite being consistently distinguishable to humans using explicit safety criteria. This suggests safety distinctions operate on semantic dimensions not captured by general-purpose pretraining.

\subsubsection{5.5 Summary}

The systematic three-hypothesis protocol reveals:
\begin{itemize}
\item \textbf{H-E1 (PASS):} Violations exist (45.6\% base-rate, p=0.0063)
\item \textbf{H-M1 (PASS):} Humans detect them consistently (κ=0.724, p=0.0076)
\item \textbf{H-M2 (FAIL):} Pretrained embeddings don't cluster them (d=0.034, p=0.797)
\end{itemize}

The causal chain breaks at embedding representation: genuine violations with consistent human detection produce no geometric structure in RoBERTa embedding space. This negative finding demonstrates that alignment evaluation via embedding clustering requires safety-specialized representations rather than standard pretrained models.
